% ! TEX root = ../m3-19-20.tex

\chapter{Șiruri de variabile aleatoare}

Pentru aspectele teoretice privitoare la aceste noțiuni, puteți consulta
\emph{Lecții de matematici speciale}, L.\ Costache.

În continuare, vom insista pe exerciții.

\section{Exerciții rezolvate}

1. Fie variabila aleatoare:
\[
  X = \begin{pmatrix}
    0.2 & 0.3 & 0.4 & 0.5 \\
    0.1 & 0.2 & 0.3 & 0.4
  \end{pmatrix}.
\]
Estimați probabilitatea $ P\left( \left\{ |X - m| < 0.2 \right\} \right) $,
unde $ m = M[X] = E[X] $ este media lui $ X $.

\emph{Soluție:} Se folosește inegalitatea lui Cebîșev, sub forma:
\[
  P \left( \left\{ | X - m| < \varepsilon \right\} \right) \geq 1 - \dfrac{\sigma^2}{\varepsilon^2},
\]
unde $ \sigma $ este dispersia, iar $ \varepsilon = 0.2 $ în acest caz (alegem).

Avem succesiv: $ M[X] = 0.4 $, deci $ (M[X])^2 = 0.16 $, iar variabila la pătrat
este:
\[
  X^2 = \begin{pmatrix}
    0.04 & 0.09 & 0.16 & 0.25 \\
    0.1 & 0.2 & 0.3 & 0.4
  \end{pmatrix}.
\]
Rezultă $ M[X^2] = 0.17 $, de unde obținem
$ \sigma^2 = D^2[X] = M[X^2] - (M[X])^2 = 0.01 $.

Din egalitatea lui Cebîșev, rezultă în fine:
\[
  P\left( \left\{ | X - m | < 0.02 \right\} \right) \geq 0.9975.
\]

\vspace{0.3cm}

2. Se aruncă 2 monede, simultan, de 800 de ori. Determinați probabilitatea
de apariție a stemei pe ambele monede de un număr de ori cuprins între
150 și 250.

\emph{Soluție:} Fie evenimentul $ A $, care arată că apare stema pe ambele
monede. Pentru o singură aruncare, avem $ p = P(A) = \dfrac{1}{4} $.

Fie acum o variabilă aleatoare cu repartiție binomială (potrivită pentru
studiul aruncării monedelor), deci de forma:
\[
  X = \begin{pmatrix}
    k \\
    C_n^k p^k q^{n-k}
  \end{pmatrix},
\]
unde $ n = 800, 1 \leq k \leq 800 $ și $ p = \dfrac{1}{4} $.

Media variabilei (folosim direct formulele cunoscute din seminariile anterioare)
este $ M[X] = np = 800 \cdot \dfrac{1}{4} = 200 $, iar dispersia se calculează cu:
\[
  \sigma^2 = np(1-p) = 150.
\]

Acum, din egalitatea lui Cebîșev:
\[
  P \left( \left\{ | X - m | < \varepsilon \right\} \right) \geq 1 - \dfrac{\sigma^2}{\varepsilon^2},
\]
unde $ m = M[X] $, trebuie să determinăm $ \varepsilon $ convenabil.

Evenimentul căutat $ |X - m| < \varepsilon $ poate fi scris în mod echivalent:
\[
  m-\varepsilon < X < \varepsilon + m \Rightarrow 200 - \varepsilon < X < \varepsilon + 200.
\]
Dar în ipoteză, avem cerința $ 150 < X < 250 $, deci se ajunge la sistemul:
\[
  \begin{cases}
    200 - \varepsilon &= 150 \\
    200 + \varepsilon &= 250
  \end{cases},
\]
sistem care are soluția $ \varepsilon = 50 $. Acum revenim în inegalitatea
lui Cebîșev și obținem:
\[
  P(150 < X < 250) \geq 1 - \dfrac{150}{50^2} \Rightarrow %
  P(150 < X < 250) \geq 0.94.
\]

\vspace{0.3cm}

3. Un zar se aruncă de 1200 de ori. Să se determine probabilitatea minimă ca
numărul de apariții ale feței cu $ i $ puncte, pentru un $ i $ fixat, să fie
între 150 și 250.

\emph{Soluție:} Fie evenimentul $ A $ care arată că apare fața cu $ i $ puncte,
pentru un $ i $ fixat, la o anume aruncare. Considerăm $ n = 1200 $ variabile
aleatoare identic repartizate, cu forma Bernoulli (potrivită pentru aruncarea
zarului), deci:
\[
  X_k = \begin{pmatrix}
    1 & 0 \\
    p & q
  \end{pmatrix},
\]
unde $ p + q = 1 $ și $ P(X_k = 1) = P(A) = p = \dfrac{1}{6} $.

Media unei asemenea variabile (din formula corespunzătoare cazului Bernoulli)
este $ M[X_k] = p $, pentru orice $ k $.

Vom folosi legea slabă a numerelor mari, în forma generală, unde frecvența
relativă de apariție a evenimentului $ A $ este $ f_{1200} = \dfrac{k}{1200} $.
Rezultă:
\[
  P \left( \left| \dfrac{k}{1200} - \dfrac{1}{6} \right| < \varepsilon \right) > 1 - \delta,
\]
unde avem $ n = \dfrac{1}{4\delta \varepsilon^2} $, care conduce la
$ \delta = \dfrac{1}{4 \cdot 1200 \cdot \varepsilon^2} $.

Pentru a determina valoarea lui $ \varepsilon $, calculăm:
\[
  \left| \dfrac{k}{1200} - \dfrac{1}{6} \right| < \varepsilon \Leftrightarrow %
  \dfrac{1}{6} - \varepsilon < \dfrac{k}{1200} < \dfrac{1}{6} + \varepsilon,
\]
de unde, în fine: $ 200 - 1200\varepsilon < k < 200 + 1200 \varepsilon $.
Deoarece cerința problemei este $ 150 < k < 250 $, se ajunge la sistemul:
\[
  \begin{cases}
    200 - 1200 \varepsilon &= 150 \\
    200 + 1200 \varepsilon &= 250
  \end{cases} \Rightarrow \varepsilon = \dfrac{1}{24}.
\]

Pentru această valoare, avem $ \delta = 0.12 $, deci $ 1 - \delta = 0.88 $.
Concluzia este că probabilitatea cerută este de 80\%.

\vspace{0.3cm}

4. Folosind teorema Moivre-Laplace, estimați probabilitatea ca din 100 de
aruncări ale unei monede, să apară banul de un număr de ori cuprins între
40 și 60.

\emph{Soluție:} Fie $ n = 100 $ variabile aleatoare repartizate identic,
binomial (model potrivit pentru aruncarea monedelor). Putem considera simplu:
\[
  X_k = \begin{pmatrix} 1 & 0 \\ p & q \end{pmatrix},
\]
unde $ P(X_k = 1) = p = 0.5 $. Media fiecăreia dintre aceste variabile este
$ M[X_k] = p $. Avem de calculat probabilitatea:
\[
  P \left( a < \sum_{k=1}^n X_k < b \right), \quad a = 40, b = 60.
\]
Aplicăm teorema Moivre-Laplace:
\[
  \ell = \lim_{n \to \infty} P \left( a < \sum_{k=1}^n X_k < b \right) = %
  \phi\left( \dfrac{b - np}{\sqrt{npq}} \right) - %
  \phi\left( \dfrac{a - np}{\sqrt{npq}} \right).
\]
Rezultă $ \ell = \phi(2) - \phi(-2) = \phi(2) - (1 - \phi(2)) $. Acum,
cu valorile din tabel pentru funcția lui Laplace, avem rezultatul cerut
ca fiind 0.9544.

\vspace{0.3cm}

5. Fie șirul de variabile aleatoare, independente $ (X_n)_{n \in \NN} $, cu
repartiția:
\[
  X_n = \begin{pmatrix}
    -n & 0 & n \\
    & & \\
    \dfrac{1}{2n^2} & 1 - \dfrac{1}{n^2} & \dfrac{1}{2n^2}
  \end{pmatrix}.
\]
Arătați că acest șir de variabile aleatoare îndeplinește condițiile legii
numerelor mari în forma Cebîșev.

\emph{Soluție:} Avem de arătat că dispersiile șirului de variabile aleatoare
sînt finite și uniform mărginite, adică există $ c \in \RR $ constantă,
astfel încît $ \sigma^2(X_k) \leq c $, pentru toți $ k $.

Dar calcule imediate arată:
\begin{align*}
  M[X_n] &= -\dfrac{n}{2n^2} + \dfrac{n}{2n^2} = 0 \\
  M[X^2_n] &= \dfrac{n^2}{2n^2} + \dfrac{n^2}{2n^2} = 1 \\
  \sigma^2 &= M[X_n^2] - M^2[X_n] = 1 - 0 = 1,
\end{align*}
care adeveresc condiția.

%%%%%%%%%%%%%%%%%%%%%%%%%%%%%%%%%%%%%%%%%%%%%%%%%%%%%%%%%%%%%%%%%%%%%%
\section{Exerciții propuse}

\indent\indent 1. Timpul mediu de răspuns al unui calculator este de 15 secunde pentru o
anumită operație, cu abaterea standard de 3 secunde.

Estimați probabilitatea ca timpul de răspuns să se abată cu mai puțin de
5 secunde față de medie.

\emph{Indicație:} Folosiți inegalitatea Cebîșev.

\vspace{0.3cm}

2. Arătați că șirul de variabile aleatore $ (X_n)_n $, independente, cu
repartiția:
\[
  X_n = \begin{pmatrix}
    -\sqrt{\ln n} & \sqrt{\ln n} \\
    0.5 & 0.5
  \end{pmatrix}
\]
urmează legea slabă a numerelor mari.

\emph{Indicație:} Verificați mai întîi dacă variabilele din șir au dispersii
finite, iar în caz negativ, verificați formularea Markov.

\vspace{0.3cm}

3. Într-o școală se află 100 de unități electrice care funcționează
independent. Probabilitatea deconectării uneia dintre ele într-un interval de
timp dat este de 5\%. Determinați probabilitatea ca în acest interval anume
să fie deconectate:
\begin{enumerate}[(a)]
\item nu mai puțin de 5 unități;
\item nu mai mult de 5 unități;
\item între 5 și 10 unități.
\end{enumerate}

\emph{Indicație:} Teorema Moivre-Laplace, pentru variabile repartizate
Bernoulli.

\vspace{0.3cm}

4. Se știe că 5\% dintre produsele unui anumit lot sînt defecte. Care
este probabilitatea ca dintr-o selecție de 200 de produse, cel puțin
10\% să fie defecte?

\emph{Indicație:} Teorema Moivre-Laplace, pentru variabile repartizate
Bernoulli.

\vspace{0.3cm}

5. O fabrică folosește 2 tipuri de aparate, $ A $ și $ B $.
După un timp, tipul $ A $ se strică și trebuie reparate cu probabilitatea de 20\%,
iar tipul $ B $, cu probabilitatea de 40\%.

Numărul total de aparate din tip $ A $ disponibile în fabrică este de 1000,
iar cele de tip $ B $ sînt în număr de 1500.

\begin{enumerate}[(a)]
\item Care este probabilitatea ca la un moment dat, numărul de aparate care
  trebuie reparate (de oricare dintre tipuri) să fie între 750 și 850?
\item Determinați, cu o probabilitate de 0.99, numărul maxim de aparate care
  ar trebui reparate la un moment dat.
\end{enumerate}

\emph{Indicație:} Teorema Moivre-Laplace, pentru variabile repartizate
Bernoulli (considerați două șiruri, corespunzătoare celor două tipuri
de aparate).


%%% Local Variables:
%%% mode: latex
%%% TeX-master: "../m3-19-20"
%%% End:
